\documentclass[12pt]{article}

\usepackage[a4paper,margin=1in]{geometry}
\usepackage{amsmath,amssymb,amsthm}

\newcommand{\submittedby}[1]{%
    \noindent\textbf{Submitted by: #1}\par
}

\begin{document}

\submittedby{Manas Singh}

\paragraph{15.} Let $\mathcal{T}_X$ denote the collection of all topologies on $X$, and let $\mathcal{T}_X/\text{hom}$ denote the set of homeomorphism classes of topologies on $X$, where $\mathcal{T}_1 \sim \mathcal{T}_2$ if and only if $(X,\mathcal{T}_1)$ and $(X,\mathcal{T}_2)$ are homeomorphic. Let $S_X$ be the group of all bijections $f: X \to X$. Each $f \in S_X$ induces the bijection $\tilde{f}: \mathcal{P}(X) \to \mathcal{P}(X)$, $\tilde{f}(A) = \{f(a) \mid a \in A\}$, and $\tilde f$ preserves arbitrary unions and intersections, with $\tilde f(\emptyset)=\emptyset$ and $\tilde f(X)=X$.

\textbf{The action is well-defined.} For $\mathcal{T} \in \mathcal{T}_X$, the collection $f\cdot\mathcal{T} = \tilde{f}(\mathcal{T}) = \{f(U) \mid U \in \mathcal{T}\}$ contains $\emptyset$ and $X$ and is closed under arbitrary unions and finite intersections, because $\tilde f$ preserves these operations. Hence $f\cdot\mathcal{T} \in \mathcal{T}_X$. Moreover $\mathrm{id}\cdot\mathcal{T}=\mathcal{T}$ and $(g\circ f)\cdot\mathcal{T} = g\cdot(f\cdot\mathcal{T})$, so this is a group action of $S_X$ on $\mathcal{T}_X$.

\textbf{Key observation.} For $f\in S_X$ and $\mathcal{T}_1,\mathcal{T}_2\in\mathcal{T}_X$,
\[
f\cdot\mathcal{T}_1=\mathcal{T}_2 \iff f:(X,\mathcal{T}_1)\to(X,\mathcal{T}_2) \text{ is a homeomorphism.}
\]
Indeed, $f$ is a bijection, and $f\cdot\mathcal{T}_1=\mathcal{T}_2$ means precisely that $U\in\mathcal{T}_1 \iff f(U)\in\mathcal{T}_2$. The forward direction ($U \in \mathcal{T}_1 \Rightarrow f(U)\in\mathcal{T}_2$) says $f^{-1}$ is continuous, and the backward direction says $f$ is continuous. So $f$ is a homeomorphism.

Conversely, every homeomorphism $f:(X,\mathcal{T}_1)\to(X,\mathcal{T}_2)$ is in particular a bijection $X\to X$, i.e.\ an element of $S_X$, and it satisfies $f\cdot\mathcal{T}_1=\mathcal{T}_2$.

\textbf{The bijection.} Define $\Phi: \mathcal{T}_X / S_X \to \mathcal{T}_X/\text{hom}$ by
\[
\Phi(\text{Orb}(\mathcal{T})) = [\mathcal{T}]_{\text{hom}}.
\]
\begin{itemize}
    \item \textbf{Well-defined and injective:} By the key observation,
    \[
    \text{Orb}(\mathcal{T}_1) = \text{Orb}(\mathcal{T}_2) \iff \exists f\in S_X:\ f\cdot\mathcal{T}_1=\mathcal{T}_2 \iff (X,\mathcal{T}_1)\cong(X,\mathcal{T}_2) \iff [\mathcal{T}_1]_{\text{hom}} = [\mathcal{T}_2]_{\text{hom}}.
    \]
    The implication ``$\Rightarrow$'' shows that $\Phi$ does not depend on the choice of representative of the orbit (well-defined), and ``$\Leftarrow$'' shows that $\Phi$ is injective.
    \item \textbf{Surjective:} Given any class $[\mathcal{T}]_{\text{hom}} \in \mathcal{T}_X/\text{hom}$, we have $\Phi(\text{Orb}(\mathcal{T})) = [\mathcal{T}]_{\text{hom}}$.
\end{itemize}
Therefore $\Phi$ is a natural bijection (it requires no choices) between the orbits of the action of $S_X$ on $\mathcal{T}_X$ and $\mathcal{T}_X/\text{hom}$.

\paragraph{16.} Let $X = \{a, b, c\}$. The group $S_3$ acts on $\mathcal{T}_X$ by $f\cdot\mathcal{T}=\{f(U)\mid U\in\mathcal{T}\}$. Every topology contains $\emptyset$ and $X$, so we list all $29$ topologies explicitly, grouped by orbit. Within each orbit, the topologies are obtained from one another by permuting $a,b,c$.

\begin{enumerate}
    \item \textbf{Indiscrete topology} (orbit size $1$):
    \begin{itemize}
        \item $\{\emptyset, X\}$
    \end{itemize}

    \item \textbf{One non-trivial open set, a singleton} (orbit size $3$):
    \begin{itemize}
        \item $\{\emptyset, \{a\}, X\}$
        \item $\{\emptyset, \{b\}, X\}$
        \item $\{\emptyset, \{c\}, X\}$
    \end{itemize}

    \item \textbf{One non-trivial open set, a two-element set} (orbit size $3$):
    \begin{itemize}
        \item $\{\emptyset, \{a,b\}, X\}$
        \item $\{\emptyset, \{b,c\}, X\}$
        \item $\{\emptyset, \{a,c\}, X\}$
    \end{itemize}

    \item \textbf{Partition topology} (orbit size $3$):
    \begin{itemize}
        \item $\{\emptyset, \{a\}, \{b,c\}, X\}$
        \item $\{\emptyset, \{b\}, \{a,c\}, X\}$
        \item $\{\emptyset, \{c\}, \{a,b\}, X\}$
    \end{itemize}

    \item \textbf{Nested chain $\{x\}\subset\{x,y\}\subset X$} (orbit size $6$, one for each ordered pair $(x,y)$ of distinct points):
    \begin{itemize}
        \item $\{\emptyset, \{a\}, \{a,b\}, X\}$
        \item $\{\emptyset, \{a\}, \{a,c\}, X\}$
        \item $\{\emptyset, \{b\}, \{a,b\}, X\}$
        \item $\{\emptyset, \{b\}, \{b,c\}, X\}$
        \item $\{\emptyset, \{c\}, \{a,c\}, X\}$
        \item $\{\emptyset, \{c\}, \{b,c\}, X\}$
    \end{itemize}

    \item \textbf{Two disjoint singletons and their union} (orbit size $3$, determined by the excluded point):
    \begin{itemize}
        \item $\{\emptyset, \{a\}, \{b\}, \{a,b\}, X\}$
        \item $\{\emptyset, \{b\}, \{c\}, \{b,c\}, X\}$
        \item $\{\emptyset, \{a\}, \{c\}, \{a,c\}, X\}$
    \end{itemize}

    \item \textbf{One singleton with both two-element sets containing it} (orbit size $3$, determined by the singleton):
    \begin{itemize}
        \item $\{\emptyset, \{a\}, \{a,b\}, \{a,c\}, X\}$
        \item $\{\emptyset, \{b\}, \{a,b\}, \{b,c\}, X\}$
        \item $\{\emptyset, \{c\}, \{a,c\}, \{b,c\}, X\}$
    \end{itemize}

    \item \textbf{Two singletons $\{x\},\{y\}$ with $\{x,y\}$ and $\{x,z\}$} (orbit size $6$, one for each ordered pair $(x,y)$ of distinct points, where $z$ is the remaining point):
    \begin{itemize}
        \item $\{\emptyset, \{a\}, \{b\}, \{a,b\}, \{a,c\}, X\}$
        \item $\{\emptyset, \{b\}, \{a\}, \{a,b\}, \{b,c\}, X\}$
        \item $\{\emptyset, \{a\}, \{c\}, \{a,c\}, \{a,b\}, X\}$
        \item $\{\emptyset, \{c\}, \{a\}, \{a,c\}, \{b,c\}, X\}$
        \item $\{\emptyset, \{b\}, \{c\}, \{b,c\}, \{a,b\}, X\}$
        \item $\{\emptyset, \{c\}, \{b\}, \{b,c\}, \{a,c\}, X\}$
    \end{itemize}

    \item \textbf{Discrete topology} (orbit size $1$):
    \begin{itemize}
        \item $\mathcal{P}(X) = \{\emptyset, \{a\}, \{b\}, \{c\}, \{a,b\}, \{b,c\}, \{a,c\}, X\}$
    \end{itemize}
\end{enumerate}

Each listed collection contains $\emptyset$ and $X$ and is closed under unions and intersections, so each is a topology. Types 1--9 are pairwise non-homeomorphic: they differ in the number of open sets ($2,3,3,4,4,5,5,6,8$) or, when these agree, in their structure (types 2 and 3 differ in the size of the non-trivial open set; types 4 and 5 differ because the open sets of type 4 are disjoint while those of type 5 are nested; types 6 and 7 differ because type 6 has two singleton open sets and type 7 has only one). Within each type, permuting the points moves any topology to any other, so each type is exactly one orbit.

Summing the orbit sizes gives
\[
1 + 3 + 3 + 3 + 6 + 3 + 3 + 6 + 1 = 29,
\]
which accounts for all $29$ topologies. Hence the number of orbits of $S_3$ on $\mathcal{T}_X$, i.e.\ the number of homeomorphism classes of topologies on a three-point set, is exactly \textbf{9}.

\paragraph{17.} Let $X$ be a set, and let $\mathcal{S}_1, \mathcal{S}_2$ be two subbases on $X$ generating topologies $\mathcal{T}_1$ and $\mathcal{T}_2$ respectively. Recall that the basis generated by a subbasis $\mathcal{S}$ is $\mathcal{B}_{\mathcal{S}} = \{ U_1\cap\cdots\cap U_n \mid n \in \mathbb{N}, U_i \in \mathcal{S} \}$, and the topology generated by $\mathcal{S}$ consists of all unions of elements of $\mathcal{B}_{\mathcal{S}}$.

\begin{enumerate}
    \item[(i)] Suppose $\mathcal{S}_1 \subseteq \mathcal{S}_2$. Let $W \in \mathcal{T}_1$. Then $W = \bigcup_{\alpha \in I} B_\alpha$, where each $B_\alpha$ is a finite intersection of elements of $\mathcal{S}_1$. Since $\mathcal{S}_1 \subseteq \mathcal{S}_2$, each set in that finite intersection belongs to $\mathcal{S}_2$, so $B_\alpha \in \mathcal{B}_{\mathcal{S}_2}$ for all $\alpha \in I$. Thus $W$ is a union of elements of $\mathcal{B}_{\mathcal{S}_2}$, i.e.\ $W \in \mathcal{T}_2$. Therefore $\mathcal{T}_1 \subseteq \mathcal{T}_2$.

    \item[(ii)] We give an explicit example where $\mathcal{T}_1 = \mathcal{T}_2$ even though $\mathcal{S}_1 \neq \mathcal{S}_2$.

    Let $X = \{1, 2, 3\}$ and define
    \[
    \mathcal{S}_1 = \bigl\{\{1\}, \{2\}, X\bigr\}, \qquad \mathcal{S}_2 = \bigl\{\{1\}, \{2\}, \{1, 2\}, X\bigr\}.
    \]
    Both cover $X$ (they contain $X$), so both are subbases. Clearly $\mathcal{S}_1 \neq \mathcal{S}_2$, since $\{1, 2\} \in \mathcal{S}_2 \setminus \mathcal{S}_1$.

    \begin{itemize}
        \item The finite intersections of elements of $\mathcal{S}_1$ are $\{1\}$, $\{2\}$, $X$ and $\{1\}\cap\{2\}=\emptyset$, so $\mathcal{B}_1 = \{\emptyset, \{1\}, \{2\}, X\}$. Taking all unions (note $\{1\}\cup\{2\}=\{1,2\}$),
        \[
        \mathcal{T}_1 = \{\emptyset, \{1\}, \{2\}, \{1, 2\}, X\}.
        \]
        \item The finite intersections of elements of $\mathcal{S}_2$ are $\{1\},\{2\},\{1,2\},X$ and $\emptyset$ (for instance $\{1\}\cap\{2\}$), and no new sets arise from intersecting these (e.g.\ $\{1,2\}\cap\{1\}=\{1\}$). So $\mathcal{B}_2 = \{\emptyset, \{1\}, \{2\}, \{1,2\}, X\}$, and since this family is already closed under unions,
        \[
        \mathcal{T}_2 = \{\emptyset, \{1\}, \{2\}, \{1, 2\}, X\}.
        \]
    \end{itemize}
    Thus $\mathcal{T}_1 = \mathcal{T}_2$ although $\mathcal{S}_1 \neq \mathcal{S}_2$.
\end{enumerate}

\end{document}
